\documentclass[11pt]{article}\usepackage[margin=1in]{geometry}\usepackage{amsmath,amsthm}\newtheorem{theorem}{Theorem}\newtheorem{remark}{Remark}
\title{Canonical Curvature Witnesses for Incoherent Collars}\author{}\date{21 July 2026}\begin{document}\maketitle
\begin{abstract}We recode every covering-radius collar failure into a deterministic gauge-invariant plaquette witness and compute its exact incidence congestion on the open $2^4$ block.\end{abstract}
For a failed link $\ell$, Stage 31 implies that the maximum curvature among incident plaquettes is greater than $\eta$. Choose the smallest-index plaquette attaining that maximum.
\begin{theorem}The selector is deterministic and gauge invariant. Every selected witness has curvature $>\eta$. On the cubical block, each plaquette can be selected by at most its four boundary links. Consequently $k$ failed links generate at least $\lceil k/4\rceil$ distinct witnesses. The capacity four is attained.\end{theorem}
\begin{proof}Curvature radii and incidence are gauge invariant. A plaquette can be incident only on its four boundary links, giving the congestion bound and the counting inequality. Setting one plaquette curvature strictly above all others makes all four boundary links select it, proving sharpness.\end{proof}
The Bianchi adjacency of the resulting plaquette objects has maximum degree 20 on the finite block. This is distinct from the shared-link bad-event graph.
\begin{remark}The theorem recodes a declared set of failed collars. It does not yet prove that an adaptive physical exposure produces a history-uniform witness law or a normalized partition activity.\end{remark}\end{document}
